\subsection{Jets}
Precise tests of QCD are challenging because the hard-scattered partons hadronize and fragment into many different hadrons.
To accurately test pQCD predictions each final-state hadron would have to be associated with a present quark in the hard scattering, which is challenging experimentally.
Luckily, for sufficiently hard processes the hard-scattered partons are sufficiently boosted in the lab frame, which means that the characteristic emissions and splittings of hadronization will happen at relatively small angles.
The consequence of this boost is that the final-state hadrons tend to form highly collimated sprays of particles known as \textit{jets}.
At leading order, jet production in $l^+l^-$ annihilation and in hard parton--parton scattering produces two outgoing partons.
In the transverse plane, momentum conservation makes these partons approximately back-to-back in azimuth.
The resulting jets are observed in back-to-back pairs called \textit{dijets}.
Jets were first observed at SLAC in $e^+e^-$ collisions~\cite{Jet:SLAC}.
Since jets are a very complicated multi-object observable, they must be defined precisely both theoretically and experimentally~\cite{Jet:QCD}.

\paragraph{Requirements for Jets}
In order to be calculable to next-to-leading order (NLO) and beyond, jets must be both infrared and collinear (IRC) safe~\cite{pdg}.
The infrared condition requires that a parton in the shower emitting a very soft ($z\to0$) parton should not change the observable.
The collinear condition requires that the splitting of one parton into two collinear partons ($\theta \to 0$) should not change the observable.
IRC safety ensures that divergences from the infrared and collinear contributions can cancel and give finite predictions.
Beyond being IRC safe, the physics community also defined some proposed standard requirements for jets at the Snowmass conference in 1990~\cite{Jet:Snowmass}.
These are as follows:
\begin{enumerate}
    \item Simple to implement in experimental analyses
    \item Simple to compute theoretically
    \item Well defined at any order of perturbation theory
    \item Yield finite cross sections at all orders of perterbation theory
    \item Yield a cross section mostly insensitive to hadronization effects
\end{enumerate}
As the field of jet physics has matured, the community has converged on a set of standard practices for defining and measuring jets.

\paragraph{Jet Kinematics} Like any physics object, the kinematics of a jet can be described by a four-vector.
The most commonly used set of variables to describe a jet is $\left(m,p_{\mathrm{T}},y,\phi\right)$, where $m$ is the jet mass, $p_{\mathrm{T}}$ is the transverse momentum, $y$ is the rapidity, and $\phi$ is the azimuthal angle around the beam axis.
The jet variables can be obtained from the energy-momentum four-vector $\left(E,p_x,p_y,p_z\right)$ through Equation~\ref{eq:jet_kinematic_mapping}\footnote{Here, we assume that the beam is oriented along the $z$ axis.}
\begin{equation}
    \label{eq:jet_kinematic_mapping}
    m = \sqrt{E^2-p^2}, \;\;
    p_{\mathrm{T}} = \sqrt{p_x^2+p_y^2}, \;\;
    y = \frac{1}{2}\ln\left(\frac{E+p_z}{E-p_z}\right), \;\;
    \phi = \operatorname{atan2}\left(p_y,p_x\right),
\end{equation}
For relativistic ($E \approx p$) jets, the rapidity can be approximated by the pseudorapidity ($\eta = -\ln{\tan{\left(\theta/2\right)}}$).
Here, $\theta$ represents the polar angle with respect to the beamline.
Working with $\phi$ and $\eta$ is convenient because they have a straightforward mapping onto a detector apparatus used for an experiment.
Values of $|\eta|$ close to zero are typically referred to as ``mid-rapidity,'' while large values of $|\eta|$ are referred to as ``forward.''
Figure~\ref{fig:atlas_dijet_event_display} shows an example of a dijet event recorded by the ATLAS detector, with the kinematics labeled.
\begin{figure}[H]
    \centering
    \includegraphics[
        width=0.9\linewidth,
        alt={ATLAS proton-proton dijet event display showing two prominent jets and their associated detector activity.}
    ]{chapters/background/sections/jets/figures/ATLASEDppDijetLabeled.pdf}
    \caption{Example of a \pxp\ dijet event recorded by the ATLAS detector. The two leading jets are visible as prominent, approximately back-to-back deposits of energy in the detector. Event display provided by the ATLAS experiment~\cite{ATLAS:UNSG-2021-48}. The labeling on the diagram is courtesy of Y. Kulinich.}
    \label{fig:atlas_dijet_event_display}
\end{figure}
\subsubsection{Jet Clustering Algorithms}
Placeholder text for this section.

\subsubsection{Jet constituents}
\paragraph{Calorimeter Jets}
\paragraph{Particle Flow Jets}
\paragraph{Track Jets}

\subsubsection{Factorization of Jet Obervables}
Placeholder text for this section.
